%----------Zusammenfassung Englisch/Abstract----------------------------------------------------------------
\addsec{Abstract}
\onehalfspacing
 {\Large The human heart is the central organ of the circulatory system and plays a vital role in supplying the body with oxygen and nutrients. Rhythmic contractions pump blood through the body, laying the foundation for all vital functions. An essential part of medical diagnostics in connection with heart function and health is the electrocardiogram (ECG). The ECG measures the electrical activity of the heart and provides valuable information about the heart rhythm and possible irregularities or diseases. As technology has advanced, the way in which ECG data can be collected and monitored has also evolved significantly. Mobile health monitoring in particular now makes it possible to record ECG data not only at specific points of time, but also in real time and over longer periods of time. Wireless transmission via Bluetooth Low Energy (BLE) enables reliable and energy-efficient transfer of data to mobile devices. This opens up new possibilities for applications that offer real-time and long-term monitoring and provide users with relevant information through graphical processing and alarm functions. The goal of this thesis is the conception and implementation of a frontend application, which realises continuous and mobile ECG monitoring.
 	Moreover this dissertation is particulary focused on implemententing  the communication and data transmission via BLE aswell as realising an 24-hour-mode, where the user is able to record his own electrocardiographic data for up to 24-hours. It also provides an broad insight into the given requirements, designprincipals, and technical challenges which accured during the implementation. In addition shows the potential for modern healthsuveillance.\par }

